%=============================================================================!
% 13.4  FRICTION AND NOISE                                                    !
%=============================================================================!
\section{Friction and noise}
\label{sec:th-langevin}

A grain of pollen in water is struck by molecules from every side; on
average the blows slow it down, a friction, and their irregularity jostles
it, a noise. Langevin's equation keeps both and nothing else of the
water. This section solves the equation for one velocity, finds the
relation between friction and noise that gives the canonical
distribution, builds the BAOAB step from it, and measures how the
friction changes diffusion and the rate of hopping.

\subsection*{Random kicks}

\begin{toolbox}[A sum of random kicks]
A velocity that receives, at each of many short intervals of length $h$,
a kick $\sigma\sqrt h\,g$, with $g$ a fresh number from the normal
distribution of mean 0 and variance 1, has after $n$ kicks the variance
$n\sigma^2h = \sigma^2t$, with $t = nh$. Each kick has the variance
$\sigma^2h$, since multiplying a quantity by a constant multiplies every
deviation from its mean by that constant, and so its variance by the
constant's square; and variances of independent kicks add
(\cref{sec:sm-probability}). The variance grows in proportion to the
time, whatever the length $h$ into which it is cut, which is why each
kick is scaled by $\sqrt h$: kicks of $\sigma h\,g$ would give the
variance $\sigma^2th$, which vanishes as $h$ shrinks. A sum of independent
Gaussian numbers is itself Gaussian (\cref{ex:th-gauss-sum}), so after
any time the velocity has a Gaussian distribution, fixed by its mean and
variance. The limit of such kicks as $h$ shrinks is the Wiener process.
\end{toolbox}

\subsection*{Langevin's equation for one velocity}

Without forces, one component of an atom's velocity under friction $\gamma$
and noise of strength $\sigma$ changes over a short interval $h$ by
\begin{equation*}
   v_{k+1} = (1 - \gamma h)\,v_k + \sigma\sqrt h\,g_k ,
\end{equation*}
the change $-\gamma v_kh$ that the drag $\dd v/\dd t = -\gamma v$ of
\cref{sec:ne-drag} makes over the interval, plus a random kick. After $n$ intervals,
covering the time $\dt = nh$, applying this $n$ times gives
\begin{equation*}
   v_n = (1 - \gamma h)^nv_0 + \sigma\sqrt h\sum_{k=0}^{n-1}(1 - \gamma
   h)^{n-1-k}g_k .
\end{equation*}
As $h$ shrinks with $\dt$ fixed, $(1 - \gamma h)^n$ tends to
$\mathrm{e}^{-\gamma\dt}$: its logarithm is $n\ln(1 - \gamma h)$, and
$\ln(1 + y)$ is $y$ plus terms of order $y^2$ for small $y$ (the Taylor
series of \cref{sec:mo-taylor}, with the slope $1/(1 + y)$ of the
logarithm, 1 at $y = 0$), so the logarithm is $-n\gamma h = -\gamma\dt$
plus $n$ errors of order $\gamma^2h^2$, together of order
$\gamma^2\dt^2/n$, which vanish as $n$ grows. The kick $g_k$ enters the sum
multiplied by $\sigma\sqrt h\,(1 - \gamma h)^{n-1-k}$, which multiplies its
variance by the square of that factor, and the variances of the
independent kicks add, so, counting them backwards with $j = n - 1 - k$,
the sum is Gaussian with the variance
\begin{equation*}
   \sigma^2h\sum_{j=0}^{n-1}(1 - \gamma h)^{2j}
   = \sigma^2h\,\frac{1 - (1 - \gamma h)^{2n}}{1 - (1 - \gamma h)^2}
   \;\to\; \sigma^2\,\frac{1 - \mathrm{e}^{-2\gamma\dt}}{2\gamma} .
\end{equation*}
The middle form is the geometric series: the sum $G =
\sum_{j=0}^{n-1}z^j$, multiplied by $z$, loses its first term, 1, and
gains $z^n$, so $zG = G - 1 + z^n$ and $G = (1 - z^n)/(1 - z)$, here with
$z = (1 - \gamma h)^2$. In
the limit $(1 - \gamma h)^{2n}$, the square of $(1 - \gamma h)^n$, tends to
$\mathrm{e}^{-2\gamma\dt}$, and $h/[1 - (1 - \gamma h)^2] = 1/(2\gamma -
\gamma^2h)$ tends to $1/2\gamma$. So over any step $\dt$, exactly,
\begin{equation}
   v(t + \dt) = c\,v(t) + \frac{\sigma}{\sqrt{2\gamma}}\sqrt{1 - c^2}\;g ,
   \qquad c = \mathrm{e}^{-\gamma\dt} ,
   \label{eq:th-ou}
\end{equation}
the Ornstein-Uhlenbeck process\cite{Uhlenbeck1930}. Over a long step $c$
is close to 0, and the velocity is Gaussian with mean zero and variance
$\sigma^2/2\gamma$, whatever it started at; \cref{ex:th-ou-variance} shows
that no other variance is kept by the step.

\subsection*{The fluctuation-dissipation relation}

The canonical distribution needs that variance to be $\kB T/m$
(\cref{eq:sm-maxwell}), so
\begin{equation}
   \sigma^2 = \frac{2\gamma\kB T}{m} :
   \label{eq:th-fdr}
\end{equation}
the noise and the friction are not independent. The friction removes
energy at a rate set by $\gamma$, the noise supplies it, and the
temperature fixes their balance. With \cref{eq:th-fdr}, \cref{eq:th-ou}
replaces $v$ by $cv + \sqrt{(1 - c^2)\kB T/m}\;g$, written $v \leftarrow cv
+ \sqrt{(1 - c^2)\kB T/m}\;g$, which takes a Maxwellian velocity to another
Maxwellian velocity: the old velocity keeps the fraction $c^2$ of the
variance and fresh noise supplies the rest, $1 - c^2$, which grows from 0
at $\gamma = 0$ towards 1 as $\gamma\dt$ grows. For argon at
\SI{135}{\kelvin}, $\sqrt{\kB T/m} = \SI{0.001676}{\angstrom\per\femto
\second}$, and with steps of \SI{10}{\femto\second}, $\gamma =
\SI{1}{\per\pico\second}$ gives $c = 0.99005$.

\subsection*{BAOAB}

With forces, Langevin's equation adds them to the drag and the noise. The
splitting of \cref{sec:in-splitting} takes each part exactly for a
while: B, a half kick by the forces; A, a half drift of the positions;
and O, the exact step \cref{eq:th-ou} of the friction and noise over the
whole step. The order B A O A B, which Leimkuhler and Matthews called
BAOAB\cite{Leimkuhler2013}, puts O between the two half drifts of
\cref{fig:th-step}, where \texttt{run} calls \texttt{Langevin.middle}:

\begin{code}
    c = math.exp(-self.friction * dt)
    sigma = np.sqrt((1 - c * c) * self._kt(self.temperature)
                    / (m * self.mv2_to_energy))
    return c * v + rng.standard_normal(v.shape) * sigma[:, None]
\end{code}

\noindent Leimkuhler and Matthews showed that this order samples
positions with unusual accuracy; for a harmonic oscillator its positions
have the exact canonical distribution at any stable step. In reduced
units, with mass, spring constant, $\kB T$ and friction all 1, so that
the period is $2\pi$, \num{4000} oscillators followed with steps of 0.5,
1.0 and 1.5 give $\ens{x^2} = 0.9986$, 0.9989 and 1.0003, against the
exact value 1, while their whole-step velocities give $\ens{v^2} =
0.9367$, 0.7503 and 0.4374, close to $1 - \dt^2/4$. A density
$\mathrm{e}^{-\tilde H/[\kB T(1 - \omega^2\dt^2/4)]}$ of the whole-step
pairs $(x, v)$, with $\tilde H$ the shadow energy of \cref{eq:in-shadow}
and $\omega = 1$, would give exactly these variances, the positions'
canonical and the velocities' lowered by that factor. A large step thus
samples the arrangements correctly while the kinetic temperature reads
low; with the steps of this book the difference is small: for the
lithium atom, steps of \SI{5}{\femto\second} against its period of
\SI{170.3}{\femto\second} give $1 - \omega^2\dt^2/4 = 0.991$, and its
measured kinetic energy under Langevin's thermostat is 0.995 of $\kB T$.

\subsection*{What the friction does to the motion}

Every velocity is pulled towards zero at the rate $\gamma$, so friction
slows diffusion as collisions did. For liquid argon (\cref{fig:th-coupling}
(a)), $\gamma = 0.1$, 1 and \SI{10}{\per\pico\second} leave 0.97, 0.71 and
0.17 of the diffusion coefficient at fixed energy, almost the same as
Andersen's thermostat at the same rates. Like collisions, the noise
changes the total momentum.

For an atom whose only bath is the thermostat, the friction also decides
how often it hops. A hop needs the barrier's energy, which only the bath
supplies, and a hop, once begun, can be turned back by the friction and
the kicks. A standard yardstick is the rate of transition-state theory,
which counts the atoms crossing the edges of a hollow's cell outwards,
from the canonical distribution\cite{Eyring1935}. In a short time $\dt$
an atom leaves through a short piece $\dd l$ of an edge if it lies within
$v_\perp\dt$ of the piece and moves outwards across it with the speed
$v_\perp$. The probability of lying in that strip, of area $v_\perp\dt\,
\dd l$, is $\mathrm{e}^{-\beta U}v_\perp\dt\,\dd l$ divided by the
integral of $\mathrm{e}^{-\beta U}$ over the cell (\cref{sec:sm-canonical}).
The velocity is independent of the position (\cref{sec:sm-maxwell}), so
$v_\perp$ is averaged on its own over the Maxwellian, the atoms that move
inwards contributing 0: with the Maxwellian's integral over all $v$,
$\sqrt{2\pi\kB T/m}$, as the divisor, the average is
$\int_0^\infty v\,\mathrm{e}^{-mv^2/2\kB T}\dd v\big/\sqrt{2\pi\kB T/m} =
\sqrt{\kB T/2\pi m}$ (\cref{ex:th-flux}). Dividing by $\dt$ and adding the
pieces round the cell gives the probability per unit time of leaving.
Each hollow's cell is a hexagon whose six edges lie halfway between
neighbouring hollows, pass through the bridges and are each $a/\sqrt3$
long, so
\begin{equation}
   k_{\mathrm{TST}} = \sqrt{\frac{\kB T}{2\pi m}}\,
   \frac{\oint\mathrm{e}^{-\beta U}\dd l}{\iint\mathrm{e}^{-\beta U}\dd A} ,
   \label{eq:th-tst}
\end{equation}
the integrals round the edges and over the cell. For lithium at
\SI{1000}{\kelvin}, $\sqrt{\kB T/2\pi m} =
\SI{0.004367}{\angstrom\per\femto\second}$, the integral round the six
edges is \SI{0.2151}{\angstrom} and that over the cell
\SI{0.7475}{\angstrom\squared}, against the cell's area of
\SI{5.2408}{\angstrom\squared}, so $k_{\mathrm{TST}} =
0.004367\times0.2151/0.7475$ per femtosecond, \SI{1.257}{\per\pico
\second}; at \SI{300}{\kelvin} it is one hop every
\SI{1.69}{\nano\second}. An atom that crosses and turns back is counted
by the theory but has not hopped, so the true rate is lower.

\Cref{fig:th-coupling}(b) counts the hops of the 1000 lithium atoms, an
atom being assigned to a hollow once it comes within \SI{0.7}{\angstrom}
of it. From $\gamma = 0.01$ to \SI{3}{\per\pico\second} the rate lies
between 0.89 and \SI{1.02}{\per\pico\second}, about 0.8 of
$k_{\mathrm{TST}}$; beyond it falls, to 0.715, 0.347 and
\SI{0.149}{\per\pico\second} at 30, 100 and \SI{300}{\per\pico\second},
more and more nearly in proportion to $1/\gamma$, as friction turns each
crossing into a slow struggle over the bridge in which the atom often
turns back. That fall is one side of Kramers' theory of
rates\cite{Kramers1940,Hanggi1990}. The theory also predicts a fall at
low friction for an atom that must gain the barrier's energy from the
bath, which the circles do not show: they start from the canonical
distribution, in which 0.194 of the atoms already have more energy than
the barrier, and on a periodic surface such atoms are never taken away,
so hops go on however weak the friction. The slow supply of energy shows
when the atoms start at the bottom of a hollow with the kinetic energy of
\SI{1000}{\kelvin} but no potential energy: at $\gamma =
\SI{0.01}{\per\pico\second}$ the thermostat supplies the missing energy so
slowly that they make only 0.228 hops per picosecond over the first
\SI{100}{\pico\second}, rising from 0.17 in the first \SI{20}{\pico
\second} to 0.35 in the last, against 1.0 at \SI{3}{\per\pico\second}.

\begin{figure}[htb]
\centering
\includegraphics{ch13_thermostats/coupling}
\caption{(a) The diffusion coefficient of liquid argon, measured from the
centre of mass, against the collision rate of Andersen's thermostat and
the friction of Langevin's, each the mean of three runs with its
standard error, against three runs at fixed energy (line and band). (b)
Hops of lithium atoms between hollows of the model surface at
\SI{1000}{\kelvin}, per atom and picosecond, against the friction:
atoms started from the canonical distribution (circles, with the
standard error over the atoms) and atoms started at the bottom of a
hollow, counted over \SI{100}{\pico\second} (squares), against the rate of
transition-state theory (dashed).}
\label{fig:th-coupling}
\end{figure}

\begin{practice}
Langevin dynamics samples the canonical ensemble reliably, even for a
single atom, and is a good thermostat for averages over arrangements and
for equilibration. A friction of \SI{1}{\per\pico\second} already slows
the diffusion of liquid argon by more than a quarter; a diffusion coefficient or
a rate measured under Langevin dynamics belongs to the friction as much
as to the system, unless the friction is small against the system's own
rates of losing momentum.
\end{practice}

\begin{keyidea}
Friction and noise together give each velocity the exact step $v
\leftarrow cv + \sqrt{(1 - c^2)\kB T/m}\,g$, $c = \mathrm{e}^{-\gamma\dt}$,
whose noise is tied to the friction by $\sigma^2 = 2\gamma\kB T/m$. BAOAB
puts it between the half drifts of velocity Verlet and samples positions
accurately even at long steps. Friction slows diffusion and, for an atom
whose bath is the thermostat, sets its rate of hopping: about 0.8 of the
rate of transition-state theory at low friction, falling more and more
nearly in proportion to $1/\gamma$ at high.
\end{keyidea}

\notebook{13}{set the friction and watch diffusion and hopping change}

\begin{selfcheck}
What is $c$ for $\gamma = \SI{10}{\per\pico\second}$ and steps of
\SI{10}{\femto\second}, and what fraction of the velocity's variance is
replaced by fresh noise at each step?
\end{selfcheck}
